\section{Related Work}

\paragraph{Personality and situated behavior.}
TRAIT extends personality inventories into scenario-based choices and evaluates
distinctiveness, consistency, and prompting effects, including agreement between
original and paraphrased items \citep{lee2025trait}. Behavioral Mode Axes
examines choices across interaction registers, evaluates open-ended generations,
and studies activation-based control \citep{liu2026bma}. Internal-feature
interventions and conversational studies further connect personality-related
behavior to changes in situations \citep{zhang2026triad,han2026contexts}.
PERSONA composes activation vectors for context-dependent personality control
\citep{feng2026persona}, while Dynamic Persona Coherence distinguishes stable
identity from appropriate state changes in role-playing \citep{qi2026dynamic}.
These studies connect consistency, context, and control. We examine that
relationship in existing AI tutors by predicting observable teaching actions
before collecting new responses and testing how shared instructions change them.

\paragraph{Tutor behavior and educational personas.}
The CIMA comparison reports characteristic action mixtures and response
complexity for three LLMs in language-learning dialogues; its generation
prompt requests both a response and action labels \citep{kucheria2025cima}.
Tutor-persona steering learns from human tutoring dialogues and evaluates
persona alignment on held-out dialogues \citep{lee2026tutorpersonas}.
These studies establish tutor-specific styles and methods for learning them.
We test whether models' past actions help predict new responses, whether the
same forecasts work when students change their wording, and how shared
instructions change both required actions and other choices.
We code visible answers after generation under controlled student cues and
teaching instructions.

\paragraph{Teaching quality and behavioral predictability.}
MRBench and MathTutorBench evaluate pedagogical capabilities
\citep{maurya2025mrbench,macina2025mathtutorbench}, while PATS studies teaching
adapted to student personality \citep{rooein2026pats}. We complement capability evaluation by asking which teaching actions can be
predicted from a model's previous responses and when changes in input make
those predictions less useful. This connects descriptions of teaching habits with evidence about what tutors
are likely to do next.
